\section{Constructing the canonical finite and real Weil actions}
\label{sec:canonical-weil}

We now construct the full representation retained in WA18.  Throughout,
the arithmetic datum is the literal $D:\texttt{D7}$, $E/F$ is its CM
extension, $\psi$ is its standard additive character of sign $-1$, and
$S,T$ are the specified Hermitian spaces.  Their ranks may be zero and
need not be even.  The given full lattices $\Lambda_S,\Lambda_T$ are
self-dual.  We retain the original infinity embeddings $\theta$, the
host ideal $\mathfrak A=\texttt{host.AA}$, and the original half
symplectic pairing $b=\texttt{symplOH}$; put $B=2b=\texttt{symplO}$.
Thus, on $X=\operatorname{Hom}_E(S,T)(\mathbb A_F)$,
\[
 H_xf(y)=\psi(b(y,x))f(y+x),\qquad
 r(g,h)x=hxg^{-1}.                                      \tag{CW.1}
\]
The theta space $\Theta$ consists of the source's continuous,
real-smooth, finite-level functions with its rational quasi-periodicity.
Statements about operators below concern this space; extending an
operator to every function on $X$ does not assert analytic properties
on that larger algebraic space.

We use ordinary finite-dimensional Fourier inversion, Haar integration,
the spectral theorem for Hermitian matrices, and elementary number-field
lattice theory.  Every infinite integral introduced below has an
explicit convergence argument.  CA3--CA4 supply the already reconstructed
finite norm Gauss sum and signed lattice-intertwiner composition.  Their
generic statements apply here without CA5's separate even-rank condition.
No full Weil representation is an input to CW1--CW22.

\paragraph{CW1: the actual finite module, integral half, and self-dual lattice.}
Let $X_f$ be the finite part of $X$, and let
\[
 A=\widehat{\mathfrak A\operatorname{Hom}_{\mathcal O_E}
                  (\Lambda_S,\Lambda_T)}\subset X_f .     \tag{CW.2}
\]
Here completion and the embedding into $X$ are the original ones; the
infinite component is zero.  WA1's integral trace-dual calculation gives
$A=A^{\perp}$ for $\gamma(x,y)=\psi(B(x,y))$.  More explicitly, local
dual bases for the two Hermitian lattices make the matrix units dual
for the trace pairing.  The factor
$\xi_0\mathfrak A\overline{\mathfrak A}$ is the host's codifferent,
so testing all integral multiples of each matrix unit gives precisely
(CW.2).  This argument works with projective lattices and does not
require a global orthonormal integral basis.

Choose $w\in\mathcal O_E$ with $w+\bar w=1$.  Such a choice is possible
because the relative trace ideal is $\mathcal O_F$: away from $2$ the
trace of $1$ is a unit locally, and at the primes over $2$ the host's
unramified quadratic algebra has surjective integral trace.  Localizing
the trace ideal proves the global assertion.  Define
\[
 W_x=\psi(-b(wx,x))H_x\quad(x\in X_f).
 \qquad W_xW_y=\gamma(wx,y)W_{x+y}.                     \tag{CW.3}
\]
To obtain the identity, expand $b(w(x+y),x+y)$, use
$b(ax,y)=b(x,\bar a y)$, skew-symmetry, and $w+\bar w=1$.
It follows that $W_xW_y=\gamma(x,y)W_yW_x$.
All these scalars have absolute value one.  On a self-dual
$\mathcal O_E$-submodule the cocycle in (CW.3) is one.

Every $r(p)A$ is self-dual and commensurable with $A$.  Indeed $r(p)$
preserves the form; its matrix and inverse are integral outside a finite
set, and at each remaining place powers of a local uniformizer bound
both images.  One rational integer clears these finitely many bounds.
The same argument shows that a sufficiently large integer multiple of
$A$ lies in any given finite open neighborhood of zero.  All ensuing
quotients of two such lattices by their intersection are finite.

\paragraph{CW2: the signed normalization is available for these lattices.}
Use the multiplicative function $\lambda:\mathbb N\to\mathbb C$
constructed in CA4.  It is nonzero at positive integers.  For each
residue field of size $q_P$, $\lambda(q_P)^2=q_P$; when conjugation
fixes $P$, its residue action is nontrivial and
$\lambda(q_P)=-\sqrt{q_P}$.  This last sign is essential at odd as
well as dyadic residue characteristic.  It is supplied by CA3's norm
Gauss sum, not chosen independently at successive lattice moves.

For clarity, the host permits a single such multiplicative function.
Its abelian Galois extension $E/\mathbb Q$ has a common residue degree
$f_p$ above a rational prime $p$, and conjugation fixes either all or
none of the primes above $p$.  Set
$\lambda(p)=\sqrt p\exp(\pi i/f_p)$ in the former case and
$\lambda(p)=\sqrt p$ in the latter, and extend by prime factorization.
The unramified CM condition makes the residue involution nontrivial
at the fixed primes.  Hence these choices have exactly the residue
values just stated.  If no rational inert prime divides $n>0$, then
$\lambda(n)=\sqrt n$.

For commensurable self-dual lattices write
\[
 n(A,C)=[C:A\cap C],\qquad
 J_{A,C}e=\lambda(n(A,C))^{-1}
       \sum_{x\in C/(A\cap C)}W_xe\quad(e\in\Theta^A).
                                                               \tag{CW.4}
\]
The sum is representative-independent because $(A\cap C)$ fixes $e$
and $C$ is isotropic.  It belongs to $\Theta^C$.  CA4 proves
$J_{C,L}J_{A,C}=J_{A,L}$ and $J_{A,A}=1$ by finite character
orthogonality and the residue norm Gauss sum.  In particular these are
isomorphisms.  The same proof, or inverse composition on a nonzero
fixed vector, gives $\lambda(n(A,C))^2=n(A,C)$ whenever that fixed
space is nonzero.  Consequently its absolute-value square is
$n(A,C)$.  This signed normalization is the one used throughout CW;
it is not replaced by positive square roots at nonsplit primes.

\paragraph{CW3: finite character projectors span the actual theta space.}
Every $f\in\Theta$ is fixed by $W_x$ for $x$ in a sufficiently small
open subgroup of $A$.  The raw finite-level condition first gives
fixing under $H_x$; intersect with the kernel of the phase in (CW.3).
By CW1 it contains $NA$ for some integer $N>0$.
Put $L=N^{-1}A$ and
$U=A\cap L^\perp=NA$.  The pairing $\gamma$ induces a perfect pairing
of the finite groups $A/U$ and $L/(A\cap L)$.  Injectivity on either
side follows from $A=A^\perp$; finite character orthogonality then
gives equality of their cardinalities and surjectivity to the dual.

For $t\in L$ form
\[
 e_t=\sum_{a\in A/U}W_aW_tf.
                                                               \tag{CW.5}
\]
The vector $W_tf$ is $U$-fixed by the commutator formula, so (CW.5)
is well-defined and fixed by $A$.  Commuting $W_{-t}$ through this
sum expresses a nonzero phase multiple of $W_{-t}e_t$ as
$\sum_{a\in A/U}\gamma(a,t)W_af$.
Sum this expression over the dual group $L/(A\cap L)$ with the
opposite sign of $t$.  Orthogonality kills every term except $a=0$,
and the remaining term is $|A/U|f$.  Division by this positive integer
proves
\[
             \Theta=\operatorname{span}_{\mathbb C}
                   \{W_xe:x\in X_f, e\in\Theta^A\}.    \tag{CW.6}
\]
Only a finite sum is used for each $f$.  The construction neither
asserts that $\Theta^A$ is one-dimensional nor requires a topological
closure of the span.

\paragraph{CW4: defining a genuine representation from the fixed vectors.}
For $p=(g,h)$ define an operator on the generators in (CW.6) by
\[
       \Omega_f(p)(W_xe)=W_{r(p)x}J_{A,r(p)A}e .          \tag{CW.7}
\]
Here is the required well-definedness argument.  An $A$-fixed vector
$e$ gives an $A$-eigenvector $W_xe$ with eigencharacter
$a\mapsto\gamma(a,x)$.  The characters of $x,x'$ are equal exactly
when $x-x'\in A$.  A finite relation between such vectors therefore
splits into its separate $X_f/A$ cosets: for finitely many distinct
characters, polynomials in the commuting $W_a$ isolate each
eigenspace.  Within one coset, move every representative to the same
$x$ using (CW.3), leaving a relation among $A$-fixed vectors.
$J_{A,r(p)A}$ is linear and preserves that relation.  The transformed
representatives have exactly the same scalar correction, since
$r(p)$ preserves $\gamma$ and commutes with $w$.  Thus (CW.7)
annihilates every relation and defines a linear operator on $\Theta$.

It intertwines $W_x$ with $W_{r(p)x}$ by direct multiplication.
Reindexing the finite sum proves
\[
 \Omega_f(p)J_{A,C}e
       =J_{r(p)A,r(p)C}\Omega_f(p)e .                  \tag{CW.8}
\]
Applying (CW.8) with $C=r(q)A$ and using CW2 proves
$\Omega_f(p)\Omega_f(q)e=J_{A,r(pq)A}e$ on $\Theta^A$,
hence the group law on all generators.  Also
$\Omega_f(1)=1$ and the inverse is $\Omega_f(p^{-1})$.
This proves the literal lattice-representation construction; it is
not an appeal to a projective representation whose multiplier remains
to be removed.

\paragraph{CW5: continuity, arbitrary commutants, and the Haar norm.}
If a linear map $L:\Theta\to\Theta$ commutes with all finite $W_x$,
then it maps $\Theta^A$ into itself and commutes with the finite sum
(CW.4).  Equations (CW.6)--(CW.7) therefore give
$L\Omega_f(p)=\Omega_f(p)L$.  No continuity assumption on $L$ is
needed for this algebraic assertion.  Real Heisenberg translations
commute with finite ones, since their cross pairings vanish; hence
they are such commutants.  The phase in (CW.3) is invariant under
$r(p)$, so covariance also holds for the untwisted $H_x$.

For each fixed $f$, $\Omega_f(p)f=f$ on a neighborhood of $p=1$.
For a generator $W_xe$, require $r(p)A=A$ and
$r(p)x-x\in NA$, with $N$ chosen to annihilate
$a\mapsto\gamma(wx,a)$ on $NA$.  These are open conditions:
use finite generators for the completed lattice and the finitely
many nonintegral local matrix entries.  Then $J_{A,r(p)A}e=e$ and
(CW.3) gives $W_{r(p)x}e=W_xe$.  Intersect finitely many such
neighborhoods for a presentation of $f$.  The group law transports
this assertion to every $p$, proving local constancy of this
unadjusted finite action on each vector.

Fix any additive Haar measure $\mu$ and any rational-lattice
fundamental domain $\mathcal F$.  The theta products
$f\overline{f'}$ are periodic and bounded on the compact quotient,
so their integrals define a finite Heisenberg-invariant Hermitian
form.  Distinct translates $W_xe,W_ye'$ of fixed vectors are
orthogonal unless $x-y\in A$: applying an $A$-element with unequal
eigencharacters proves the assertion.  Consequently (CW.4) multiplies
the pairing of two fixed vectors by
$n(A,C)/|\lambda(n(A,C))|^2=1$.  Covariance, (CW.6), and the same
orthogonality prove the equality for arbitrary pairs.  Thus
\[
 \int_{\mathcal F}|\Omega_f(p)f|^2\,d\mu
                     =\int_{\mathcal F}|f|^2\,d\mu .   \tag{CW.9}
\]
This retains the given Haar measure; no unrecorded volume rescaling
has been made.  A zero fixed space causes no division problem, since
all its assertions are then zero and (CW.6) handles the whole space.

\paragraph{CW6: the exact positive split scalar and nonsplit fixing.}
Let $e$ be a split-place idempotent, $P x=ex$, $Qx=\bar e x$, and
suppose $r(p)$ is the identity off $P+Q$.  Both projections preserve
$A$, satisfy $PQ=QP=0$, and are adjoint for $\gamma$.
Let $\ell$ be any linear functional invariant under $H_{Qx}$ on
theta functions.  The phase in (CW.3) is one on this half-space,
so $\ell$ is also invariant under $W_{Qx}$.  With representatives
$b_i$ for $r(p)A/(A\cap r(p)A)$ put
\[
 n_e=\#\{i:Pb_i\in A\}>0.
 \qquad
 \ell(\Omega_f(p)f)
       ={n_e\over\lambda(n(A,r(p)A))}\ell(f).          \tag{CW.10}
\]
To prove it first take $f$ fixed by $A$.  Remove the $Qb_i$ factor
using $\ell$-invariance, and the component off $P+Q$ using
$A$-fixing.  The remaining term is $\ell(W_{Pb_i}f)$.
It equals $\ell(f)$ when $Pb_i\in A$.  Otherwise self-duality
provides $a\in A$ with $\gamma(Pb_i,a)\ne1$; its $Qa$ component
forces that term to vanish.  The zero quotient class gives $n_e>0$.
For $W_xf$, decompose $x$ into $Px,Qx,(1-P-Q)x$ and apply the
same argument to $\ell W_{(1-P-Q)x}$; covariance gives the identical
constant.  CW3 completes the proof.

The index in (CW.10) is supported at this split pair.  A rational
prime inert in the host cannot occur in that support: conjugation
has the same splitting behavior at every prime above it, because
the host extension is abelian.  CW2 therefore makes its denominator
the positive square root of the index.  Thus the multiplier in
(CW.10) is strictly positive, independent of both $f$ and $\ell$.

At a nonsplit finite place $v$, let $P_v$ be its idempotent.
The host makes every dyadic place split, so $P_v=2b_v$ for an
integral finite adele $b_v$.  If the pair $p$ is supported at $v$
and preserves $P_vA$, it preserves all of $A$, since it is the
identity on $(1-P_v)A$.  A vector fixed by the raw $H_x$ for
$x\in P_vA$ is fixed by $W_x$ there: write $x=2b_va$ and use
isotropy of $A$ to obtain $\psi(-b(wx,x))=1$.  The elementary
local-projector argument now gives $\Omega_f(p)f=f$ on these
vectors.  One proof is to apply the finite projectors of CW3 to
such an $f$: components with nontrivial character on $P_vA$ vanish;
in each remaining component the $P_v$ translation lies in $A$ by
self-duality and can be absorbed into the fixed vector.  The
remaining translation is outside $v$, where $r(p)$ is the identity.
This is the literal local-vacuum condition, with the same completed
Hom lattice, rather than merely invariance under a smaller subgroup.

\paragraph{CW7: the crossed determinant character constructs IsFinWeil.}
Let $\chi_S,\chi_T$ be the given unitary Hecke characters with
the exact restrictions of ranks $\dim_E S,\dim_E T$.  Choose one
member of every split pair and let $\epsilon$ be its finite adele
idempotent, zero at all nonsplit and infinite places.  The choices
form an adele because all entries are integral zeros or ones.  For
a unit $u$ write
\[
 \operatorname{cut}_\epsilon(u)=\epsilon u+1-\epsilon,
 \qquad
 z(g,h)=\chi_T(\operatorname{cut}_\epsilon\det g)^{-1}
         \chi_S(\operatorname{cut}_\epsilon\det h).
                                                               \tag{CW.11}
\]
The cut is a continuous group homomorphism, its inverse being the
same cut of $u^{-1}$.  The rank restriction implies that each
$\chi$ kills $u\bar u$: its restriction to $F$ is a power of the
quadratic norm character, which is one on norms.  Thus selecting
the conjugate member of a supported split pair gives the same
character value, because a unitary determinant satisfies
$\det g\,\overline{\det g}=1$.  At a real or nonsplit-supported
element the cut determinant is one.

Set $\omega_f(p)=z(p)\Omega_f(p)$ on $\Theta$.  Equations
(CW.7)--(CW.11) prove its group law, pointwise continuity, norm
preservation, real triviality, covariance, and nonsplit vacuum
fixing.  At a split pair (CW.10) becomes the exact positive
multiple of
$\chi_T(x)^{-1}\chi_S(y)=\texttt{splitFacO}(\chi_S,\chi_T,x,y)$,
where $x,y$ are the original determinant lifts.  The characters are
deliberately crossed.

To obtain the literal endomorphism of the whole function space,
choose an algebraic complement $\mathcal C$ to $\Theta$ and act as
the identity on $\mathcal C$.  Equivalently, for a projection
$\pi$ onto $\Theta$, the extension is
$1+\iota(\omega_f(p)-1)\pi$.  Its multiplication law follows from
$\pi\iota=1$.  All required assertions are on $\Theta$ and are
unchanged.  This constructs every clause of \texttt{IsFinWeil}.

\paragraph{CW8: Gaussian projectors and their actual smooth domain.}
Let $V$ be a finite-dimensional Euclidean real space with an
orthogonal complex structure $J$, and let $\rho$ be a strongly
continuous isometric Heisenberg action on a complex Banach space
$\mathcal B$:
\[
 \rho(u)\rho(v)=e^{\pi i\langle u,Jv\rangle}\rho(u+v).
 \quad G(u)=e^{-\pi\|u\|^2/2},\quad
 P f=\int_VG(u)\rho(u)f\,du .                           \tag{CW.12}
\]
Lebesgue measure is the Euclidean normalization.  The integral is
absolutely convergent since $\|\rho(u)f\|=\|f\|$.
The one-variable Gaussian Fourier formula, applied in orthogonal
coordinates, gives $P^2=P$ and
$P\rho(v)P=G(v)P$.  For example the coefficient of $\rho(x)$
in $P^2$ is
$\int G(u)G(x-u)e^{\pi i\langle u,Jx\rangle}\,du=G(x)$;
completion of the square gives the last equality with precisely
the normalization in (CW.12).

Put $\mathcal B^\infty=\{f:v\mapsto\rho(v)f\text{ is }C^\infty
\text{ in Banach norm}\}$.  For $f$ in this space,
$F_f(u)=G(u)\rho(u)f$ is a Banach-valued Schwartz function.
Indeed
\[
 F_f(t+v)=G(v)\rho(v)
     e^{-\pi\langle v,t\rangle-\pi i\langle v,Jt\rangle}F_f(t).
                                                               \tag{CW.13}
\]
Differentiating at $t=0$ bounds each derivative at $v$ by a
polynomial in $\|v\|$ times $G(v)$ and finitely many derivative
norms at zero.  Every additional polynomial factor is bounded
against the Gaussian.  Fourier inversion is therefore legitimate
in Banach norm.  The Heisenberg law gives
\[
 P\rho(-v)f=\rho(-v)\widehat F_f(Jv),\qquad
       f=\int_V\rho(v)P\rho(-v)f\,dv .                 \tag{CW.14}
\]
In particular all polynomial moments of
$\|P\rho(-v)f\|$ are integrable.

Conversely these moment bounds imply $f\in\mathcal B^\infty$:
(CW.14)'s first identity makes all polynomial moments of
$\widehat F_f$ integrable, so differentiation under Fourier
inversion makes $F_f$ smooth.  Divide by the nonvanishing smooth
Gaussian.  This equivalence is the source's \texttt{Sm}/\texttt{Good}
criterion.  It supplies both the convergence and the smooth domain
for the operators in the next paragraphs.
\paragraph{CW9: evaluating the two Gaussian scalars and their phases.}
Choose complex Euclidean coordinates in which $Jz=is z$, where
$s=s^*=s^{-1}$.  For $g^*sg=s$ put
\[
 p={1+s\over2},\quad q=1-p,\quad
 c(g)=\det(pg+q),\quad
 L(g)=\int G(u)G(gu)\,du .                             \tag{CW.15}
\]
For two such matrices set
\[
 M(g_1,g_2)=\int\!\int G(u)G(v)
       e^{\pi i\langle g_1u,Jg_2v\rangle}G(g_1u+g_2v)\,du\,dv.
                                                               \tag{CW.16}
\]
The exact evaluations are
\[
 \begin{split}
 c(g^{-1})&=\overline{c(g)},\qquad
 L(g)&=[c(g)c(g^{-1})]^{-1},\\
 M(g_1,g_2)&=[c(g_1^{-1})c(g_2)c(g_2^{-1}g_1)]^{-1}.
 \end{split}                                                   \tag{CW.17}
\]
Here $c(g)\ne0$.  These formulas determine a phase, as well as
an absolute value.

We give the matrix calculation behind them.  The exponents in
(CW.15)--(CW.16) are $-\pi z^*Az$ and $-\pi(z,z')^*N(z,z')$,
where
\[
 A={1+g^*g\over2},\qquad
 N=\begin{pmatrix}
 (1+g_1^*g_1)/2&g_1^*pg_2\\
 g_2^*qg_1&(1+g_2^*g_2)/2
 \end{pmatrix}.                                                \tag{CW.18}
\]
$A$ is positive definite and
$N+N^*=1+[g_1\ g_2]^*[g_1\ g_2]$ is positive definite.
For any complex matrix with positive Hermitian part,
\[
              \int_{\mathbb C^n}e^{-\pi z^*Az}\,dz=(\det A)^{-1}.
                                                               \tag{CW.19}
\]
For positive Hermitian $A$, diagonalize and use the elementary
complex one-dimensional Gaussian integral.  Both sides are
holomorphic in the matrix entries on the connected positive-real-part
domain; local Gaussian domination permits differentiation under the
integral.  Agreement on its open positive Hermitian real slice
implies agreement in a complex neighborhood by the power series
identity theorem, and then throughout that domain.  This proves
(CW.19), including absolute integrability, without choosing a square
root of a determinant.

For completeness the determinant reductions in (CW.18) use only
$p^2=p$, $pq=0$, and $g^*=sg^{-1}s$.  If $h=g^{-1}$, then
\[
 A=sh(pgp-qgq),\qquad
 \det(pgp-qgq)=\det(s)c(g)c(h)\det(g).
\]
Thus $\det A=c(g)c(h)$; taking adjoints gives
$c(h)=\overline{c(g)}$.  For the second matrix factor out
$\operatorname{diag}(sg_1^{-1},sg_2^{-1})$.  The remaining matrix is
\[
 K=\begin{pmatrix}pg_1p-qg_1q&pg_2\\-qg_1&pg_2p-qg_2q\end{pmatrix}.
\]
Conjugate by $\begin{pmatrix}p&q\\q&p\end{pmatrix}$, whose square
is one.  This gives an upper triangular block matrix with diagonal
blocks $sg_2(g_2^{-1}g_1p+q)$ and $pg_2p-qg_1q$.
Taking their determinants gives
$\det N=c(g_1^{-1})c(g_2)c(g_2^{-1}g_1)$.
Equation (CW.19) now gives all of (CW.17).

\paragraph{CW10: an actual oscillator action, rather than a projective one.}
For a group acting through the matrices in CW9, let $P_g$ be
(CW.12)'s projector for $u\mapsto\rho(gu)$ and define on
$\mathcal B^\infty$
\[
 T_g=c(g)P_gP,\qquad
 \Omega_\infty(g)f=\int_V\rho(gv)T_g\rho(-v)f\,dv .      \tag{CW.20}
\]
This integral converges absolutely because
$\|T_g\rho(-v)f\|\le |c(g)|\|P_g\|\|P\rho(-v)f\|$.
There is an exact coherent-coefficient identity, stronger than a
kernel estimate:
\[
 P_g\rho(-gv)\Omega_\infty(g)f=T_g\rho(-v)f.
\]
Indeed insert (CW.20), use $P_gT_g=T_g$, and apply
$P_g\rho(g(w-v))P_g=G(w-v)P_g$ with the Heisenberg phase.
The resulting integral is also obtained by applying
$T_gP\rho(-v)$ to the resolution (CW.14), since $T_gP=T_g$.
Both operations commute with the integral by the preceding absolute
bound.  For every integer $k\ge0$ this gives the explicit estimate
\[
 \int\|v\|^k\|P_g\rho(-gv)\Omega_\infty(g)f\|\,dv
 \le \|T_g\|\int\|v\|^k\|P\rho(-v)f\|\,dv<\infty.
\]
Thus the result satisfies CW8's moment criterion for $\rho\circ g$.
It is smooth for that action, and composition with $g^{-1}$ proves
smoothness for $\rho$.  All later linearity and covariance operations
are therefore on the genuine smooth domain.

Here are the identities that verify the group law.  Completing
the square in the absolutely convergent double and triple
Gaussian integrals gives
\[
 PP_gP=L(g)P,\qquad
 PP_{g_1}P_{g_2}P=M(g_1,g_2)P.                         \tag{CW.21}
\]
Also (CW.14), followed by translation of the integration variable,
gives
\[
 P_g\Omega_\infty(g)f=T_gf,\quad
 \Omega_\infty(g)\rho(v)f=\rho(gv)\Omega_\infty(g)f,
 \quad \Omega_\infty(g)P_h=P_{gh}\Omega_\infty(g).       \tag{CW.22}
\]
CW9 says exactly
\[
 |c(g)|^2L(g)=1,\qquad
 c(g)c(h)M(gh,g)=c(gh)L(gh).                           \tag{CW.23}
\]
Applying $P$ and using (CW.21)--(CW.23) shows that the difference
between $c(g)c(h)P_{gh}P_gP$ and $c(gh)P_{gh}P$ is killed by $P$.
It is in the range of $P_{gh}$; applying
$P_{gh}PP_{gh}=L((gh)^{-1})P_{gh}$ shows that it is zero.
Use this equality in (CW.22), and test after every translated
$P_{gh}$.  Formula (CW.14) says these tests separate smooth vectors.
The result is
\[
 \Omega_\infty(g)\Omega_\infty(h)=\Omega_\infty(gh),
                   \qquad\Omega_\infty(1)=1.           \tag{CW.24}
\]

Every bounded linear map commuting with $\rho$ commutes with
$P_g,T_g$ and (CW.20).  Every bounded linear map into a Hilbert space
whose norm is Heisenberg invariant is preserved by (CW.20).
For the latter assertion $P,P_g$ are self-adjoint for the resulting
sesquilinear form; (CW.21) and (CW.23) show $T_g$ is isometric on
the range of $P$.  The resolution (CW.14) gives
\[
 \langle f,f'\rangle
   =\int\langle P\rho(-v)f,P\rho(-v)f'\rangle\,dv.
\]
Use (CW.22) and the analogous resolution for $\rho\circ g$ to
deduce equality of the inner products after $\Omega_\infty(g)$.
All double integrals here are justified by CW8's moment bounds
and Gaussian domination.  Finally, (CW.20) depends continuously
on $g$ for each fixed smooth $f$: on a compact neighborhood of
$g_0$, the projector norms are uniformly bounded, and the same
integrable function $C\|P\rho(-v)f\|$ dominates the integrands.

\paragraph{CW11: the Cartan-vacuum determinant and change of coordinates.}
The vacuum normalization in CW10 can be checked directly.  If
\[
 d\rho(z)f+i\,d\rho(Jz)f=0\quad\text{for every }z,
                                                               \tag{CW.25}
\]
then $Pf=f$.  Indeed, for each $u$, the function
$\zeta\mapsto e^{\pi|\zeta|^2\|u\|^2/2}
 \rho((\Re\zeta)u+(\Im\zeta)Ju)f$
is holomorphic: differentiating and using (CW.25) is the
Cauchy--Riemann equation.  Its circle mean is its value $f$ at
zero.  Therefore the circle mean of the unweighted orbit on
$|\zeta|=1$ is $G(u)f$.  Average the integral (CW.12) over those
rotations; Fubini is bounded by $G(u)\|f\|$.  It becomes
$\bigl(\int e^{-\pi\|u\|^2}\,du\bigr)f=f$.

If $g$ is Euclidean unitary, $P_g=P$.  Equations (CW.20)--(CW.22)
then give $\Omega_\infty(g)f=c(g)f$ on this vacuum space.
For a different Cartan involution $s'$ with $Hs'$ positive,
choose an $H$-unitary change of coordinates taking $s$ to $s'$.
This exists by choosing orthonormal bases on the positive and
negative eigenspaces for the two Cartan decompositions.  Conjugate
the preceding assertion and use (CW.24).  Since determinant is
unchanged by conjugation, an element commuting with $s'$ acts
on its vacuum space by
\[
             d_P(s',g)=\det\left({1+s'\over2}g
                                      +{1-s'\over2}\right).
                                                               \tag{CW.26}
\]
It is this determinant itself, not its reciprocal.

For an arbitrary nonsingular Hermitian matrix $H$, choose an
invertible complex matrix reducing it by congruence to a signature
involution.  Pull $\rho$, the real coordinates, the group matrices
and the vacuum equations through that same linear equivalence.
CW8--CW10 are preserved, and (CW.26) is unchanged because its
matrix is conjugated.  Thus no definiteness of $S$ or $T$ is
needed for the real construction.

\paragraph{CW12: theta smoothness really is Banach smoothness here.}
Let $\mathcal B$ be the closed subspace of bounded continuous
functions on $X$ satisfying the original rational quasi-periodicity,
with supremum norm.  Compactness of $X/X(E)$ bounds every continuous
theta function; the quasi-periodic factors have modulus one.
The action $H_x$ is isometric and strongly continuous on this space.
For the latter assertion use uniform continuity on a compact set
meeting all rational cosets and the continuous phase in (CW.1).
Point evaluations are bounded and norming.

Use the actual inverse infinity chart of $\theta$ to write the
real Heisenberg action as $\rho(v)$.  A theta function has pointwise
directional derivatives, and each derivative is again a theta
function: differentiate the quasi-periodic formula and finite-level
identities, and then the defining real smoothness condition.
Denote the derivative at zero by $D_fz$.  Evaluations show this
is real linear in $z$; its finite-dimensional domain makes it
continuous.  The Heisenberg law gives the candidate Banach derivative
\[
 d(\rho(\cdot)f)_v(z)=\rho(v)D_fz
                         -\pi i\langle v,Jz\rangle\rho(v)f.
                                                               \tag{CW.27}
\]
It is continuous in $v$ in operator norm: this follows from strong
continuity on finitely many basis derivatives $D_fe_i$ and $f$.
Apply the scalar mean-value inequality to every evaluation of the
remainder, uniformly bounded by the modulus of continuity of
(CW.27).  Taking the supremum proves the Fr\'echet derivative in
Banach norm.  Iterating on the derivatives proves
$f\in\mathcal B^\infty$.  This is the missing uniform argument;
pointwise differentiability by itself would not suffice.

Conversely a Banach-smooth vector in $\mathcal B$ that is fixed at
some finite open level is a literal theta function.  Evaluate its
smooth real orbit and cancel the explicit smooth phase in (CW.1)
to get the source's real-line smoothness.  Quasi-periodicity and
continuity were already part of $\mathcal B$.  Thus the operators
of CW10 may be applied to precisely the functions needed here.

\paragraph{CW13: constructing IsInfWeil on the original theta functions.}
Vectorizing the actual real Hom matrices gives the complex index
set $((\operatorname{Fin}\dim T\times\operatorname{Fin}\dim S)
\times\mathrm{Tau})$.  In these coordinates the matrix of the
unitary pair is
\[
 m(g,h)=\operatorname{diag}_\tau
   \bigl(h_\tau\otimes(g_\tau^{-1})^{\mathsf t}\bigr).
                                                               \tag{CW.28}
\]
The real alternating form is the imaginary part of the original
nondegenerate Hermitian Hom form, including its prescribed host
scale and the sign of $\psi$.  Expanding the trace pairing proves
$m^*Hm=H$.  Apply CW10--CW11 to these matrices and to CW12's
Banach space.

Finite Heisenberg operators commute with the real action: their
cross pairings are zero, so they commute with every real $\rho(v)$
and hence with the Gaussian operator of CW10.  A finite open level
fixing $f$ therefore fixes its image.  CW12 makes that image a theta
function.  Restrict the resulting linear action to $\Theta$, then
extend by the identity on an algebraic complement, exactly as in
CW7.  Call the result $\omega_\infty$.

It is trivial on the finite subgroup, since (CW.28) is then the
identity.  Covariance for real shifts follows from (CW.22), and
for finite shifts from the commutant assertion.  Splitting an
arbitrary shift into its two orthogonal parts proves full
Heisenberg covariance for a real group element.  Pointwise
continuity follows from CW10 followed by evaluation.  For every
Haar measure and rational fundamental domain, the map
$\mathcal B\to L^2(\mathcal F,\mu)$ is bounded because
$\mu(\mathcal F)<\infty$; its norm is Heisenberg invariant by
translation on the compact quotient.  CW10 therefore proves
the exact source norm equality.

Finally the Cartan involution on the Hom coordinates is
$s=\operatorname{diag}_\tau(s_{T,\tau}\otimes
s_{S,\tau}^{\mathsf t})$.  The source's Fock equations are
exactly (CW.25), as evaluation of (CW.27) at zero shows.
For a pair commuting with these Cartans, (CW.26) equals
\[
 \prod_\tau\det\left[
 {1+s_{T,\tau}\otimes s_{S,\tau}^{\mathsf t}\over2}
       (h_\tau\otimes(g_\tau^{-1})^{\mathsf t})
 +{1-s_{T,\tau}\otimes s_{S,\tau}^{\mathsf t}\over2}\right].
                                                               \tag{CW.29}
\]
This is literally \texttt{VacCharPlus}.  All clauses of
\texttt{IsInfWeil}, on the original functions and original charts,
have now been constructed.

\paragraph{CW14: commutation and continuity of the combined action.}
For our actual finite action,
$\omega_f(p)=z(p)\Omega_f(p)$, CW5 immediately implies
\[
            \omega_f(p)\omega_\infty(q)f
             =\omega_\infty(q)\omega_f(p)f\quad(f\in\Theta),
                                                               \tag{CW.30}
\]
because the real action commutes with every finite $W_x$.
Hence $\omega(p)=\omega_f(p)\omega_\infty(p)$ is a genuine
representation on $\Theta$.  Extend this product representation
itself to all functions by one complement.  The arbitrary
extensions of the two factors need not commute outside $\Theta$.

There is a useful direct proof of the required pointwise continuity
which avoids an unjustified joint continuity assertion.  Near
$p_0$, write $p=p_0q$.  For this fixed $f$, CW5 gives
$\Omega_f(q)f=f$ when $q$ is sufficiently near $1$.  Using (CW.30)
and commuting only fixed operators gives
\[
 \omega(p)f=z(q)\,\omega_\infty(p_0q)\omega_f(p_0)f.
                                                               \tag{CW.31}
\]
The right side is pointwise continuous by CW7 and CW13.
Each factor preserves the exact Haar norm, so their product does
as well.  The real and finite subgroup restrictions are exactly
$\omega_\infty$ and $\omega_f$ respectively.  Splitting both a
group element and a Heisenberg shift into real and finite parts
proves full covariance under $r(p)$.  At this stage rational
automorphy remains to be proved; it is not a consequence of
(CW.30) or of the norm equality.

\paragraph{CW15: the scalar lemma with its essential norm hypothesis.}
Suppose $M$ preserves $\Theta$, commutes with every $H_x$, and
preserves its Haar $L^2$ norm.  Then on $\Theta$ it is multiplication
by one scalar of absolute value one.  Polarization first proves
preservation of the integrated pairings.  For $a\in X(E)$ the
quasi-periodicity gives
\[
                    H_af(x)=\psi(B(x,a))f(x).           \tag{CW.32}
\]
Consequently every Fourier coefficient of the continuous periodic
function $(Mf)\overline{Mf'}-f\overline{f'}$ is zero.
The characters in (CW.32) separate points of $X/X(E)$: the
annihilator of the rational lattice for the standard trace character
is itself.  Here is an elementary check of the underlying assertion.
For $\mathbb Q$, subtract a rational number to make every finite
component integral.  Pairing with the integers then makes the real
component an integer; subtract it.  Pairing successively with
$p^{-k}$ forces the remaining $p$-component into every $p^k\mathbb Z_p$,
hence to zero.  A rational basis and its trace-dual basis reduce a
number field, and then the Hom space, to this calculation.
Character orthogonality and the density of their
self-conjugate algebra in continuous functions on this compact
quotient show that the function is zero.  One can use uniform
approximation followed by integration against its complex conjugate;
this needs no assumption about the spectrum of $M$.

Apply this identity to $f$ and $H_{y-x}f'$ and cancel the nonzero
phase.  It follows that
$(Mf)(x)\overline{(Mf')(y)}=f(x)\overline{f'(y)}$ for all $x,y$.
If $f_0(y)\ne0$, this determines the same scalar for every $f,x$,
and the case $f=f'=f_0,x=y$ gives absolute value one.  If the theta
space is zero, take the scalar $1$.

For a rational unitary pair $p$ let
$R_pf(x)=f(r(p^{-1})x)$.  It preserves quasi-periodicity, finite
level and real smoothness.  Its rational linear change of variables
preserves the lattice quotient and Haar measure: the adelic
modulus of its rational determinant is one by the product formula.
It has the same covariance as $\omega(p)$.  Apply the scalar lemma
to $\omega(p)R_{p^{-1}}$.  Thus
\[
                     \omega(p)f=c_pR_pf,
                              \qquad |c_p|=1.          \tag{CW.33}
\]
The next calculations determine $c_p$; no assertion that it is
algebraic or equal to one has been used.

\paragraph{CW16: the finite Hermitian Gauss sum with its signed value.}
Let $M$ be a finite $\mathcal O_E$-module with a nondegenerate,
symmetric bicharacter $\beta$, satisfying
$\beta(ax,y)=\beta(x,\bar a y)$.  Put
$Q(x)=\beta(wx,x)$.  The integral half in CW1 gives
\[
 Q(x+y)=Q(x)Q(y)\beta(x,y),\qquad
       \sum_{x\in M}Q(x)=\lambda(|M|),\quad
                        \lambda(|M|)^2=|M|.           \tag{CW.34}
\]
Here the same $\lambda$ as in CW2 is used.

We prove the two sum assertions together by induction on $|M|$.
The zero module gives $1$.  Otherwise choose a simple submodule
$D\simeq\mathcal O_E/P$, of order $q$.
If $\beta|_{D\times D}=1$, then $Q|_D=1$.
Summing $Q(x+d)$ over $d\in D$ kills $x\notin D^\perp$.
On $D^\perp$ the function descends to the nondegenerate module
$D^\perp/D$, whose order is $|M|/q^2$ by perfect character duality.
Thus the sum is $q$ times its Gauss sum.  Induction and
$\lambda(q)^2=q$ give both assertions in (CW.34).

If the restriction is nontrivial, its radical is an
$\mathcal O_E$-submodule of the simple module $D$, hence zero.
Perfect duality gives $M=D\oplus D^\perp$.  Adjointness forces
$\bar P=P$: otherwise $P+\bar P=\mathcal O_E$ would make the
restriction zero.  The residue involution is nontrivial by the
unramified host condition.  On $D$, the quadratic function is
$a\mapsto\chi(w a\bar a)$ for a nontrivial additive character
with the required conjugation invariance.  CA3's explicit norm
sum is $-\sqrt q=\lambda(q)$, also in characteristic two.
The sums on the two orthogonal factors multiply, proving (CW.34).
This also proves the version for a finite quotient $B/C$ when
$C$ is exactly the radical of $\beta|_B$: the formula descends,
and a transversal sums it once.  Thus no semisimplicity of the
whole finite module, and no residue-characteristic restriction,
is hidden in the argument.

\paragraph{CW17: paired rational lattice indices and the actual finite sum.}
Let $V=\operatorname{Hom}_E(S,T)$ and
$L=\{x\in V:x_f\in A\}$, the source's \texttt{latQ}.
It is a full finitely generated $\mathbb Z$-lattice.  Its finite
closure is $A$: choosing local bases, finite approximation with
integral denominators gives density, and a common integer bounds
the finitely many nonstandard local lattices.  Consequently
$L$ is self-dual for
$\gamma_Q(x,y)=\gamma(x_f,y_f)$.  The relative index of two such
rational lattices equals the relative index of their completions,
since a finite lattice quotient is unchanged by finite completion.

Fix a rational unitary pair $p=(g,h)$.  Write $a$ for the rational
matrix of $g$ and $b_0$ for that of $h^{-1}$.  Throughout this
calculation
\[
 u(x)=b_0xa=r(p^{-1})x,\quad
 U=b_0\otimes a^{\mathsf t},\quad
 m=\dim_E V,\quad d=\det(U-1),\quad\Delta=\det U.
                                                               \tag{CW.35}
\]
Assume $d\ne0$; $\Delta\ne0$ follows from unitarity.
Put
\[
 B_1=(u-1)^{-1}L,\quad B_2=u^{-1}L,\quad
 C=L\cap B_1=L\cap B_2,\quad n=[B_1:C],\quad i=[B_2:C].
                                                               \tag{CW.36}
\]
The equality of intersections is simply
$ux-x\in L\iff ux\in L$ when $x\in L$.
Moreover $i=[r(p)A:A\cap r(p)A]$, with the orientation used in
(CW.4).  Choose an integer $N>0$ with $NuL\subset L$.
Then $NL\subset C$; let $c_0=[C:NL]>0$.
Applying $u-1$ and $u$ to the corresponding lattice quotients
and taking determinants over $\mathbb Q$ gives
\[
 c_0n=|N_{E/\mathbb Q}(N^m d)|,\qquad
 c_0i=|N_{E/\mathbb Q}(N^m\Delta)|.                   \tag{CW.37}
\]
Both equalities have the same $N,c_0,m$.  This common denominator
is needed for the later phase cancellation.

For $x,y\in B_1$ define
\[
 \beta(y,x)=\gamma_Q(\bar w(uy-y),x)
             \gamma_Q(w(uy-y),ux)
          =\gamma_Q(uy-y,wx+\bar w ux).
\]
Expansion, skew-symmetry and $\gamma_Q(ux,uy)=\gamma_Q(x,y)$
show that this bicharacter is symmetric and conjugate-adjoint.
Its radical is exactly $C$.  Indeed
$wx+\bar w ux=x+\bar w(ux-x)$; since $u-1$ maps $B_1$
onto $L$, orthogonality against all of $B_1$ is equivalent,
by self-duality of $L$, to $x+\bar w(ux-x)\in L$,
and hence to $x\in L$.  Apply CW16 to $B_1/C$.
Using $w+\bar w=1$ and $u$-invariance once more simplifies its
quadratic function to
$\beta(wt,t)=\gamma_Q(\bar w(ut-t),t)$.
Thus, for any transversal $t_1,\ldots,t_n$,
\[
       \sum_{k=1}^n\gamma_Q(\bar w(ut_k-t_k),t_k)
                       =\lambda(n),\qquad\lambda(n)^2=n.
                                                               \tag{CW.38}
\]
All these are integral quotient cardinalities, not dimensions of
rational spans.

\paragraph{CW18: the finite and infinite phases of the same eta character.}
Now take the supplied character $\eta$ satisfying the literal
\texttt{IsEtaTilde}: it is Hecke, unramified at every finite place,
has \texttt{ResEta} with exponent $1$, and infinity exponent
$1/2$ in the source convention.  The last convention means the
phase exponent is $2(1/2)=1$.

The local signs can be obtained without a new class-field theorem.
HF9 constructs the character on $\mathbb A_F^\times$ having
local factors $(-1)^{\operatorname{ord}_v}$ at inert primes,
$1$ at split primes, and the sign at each real place.
It is continuous; HF9 proves that it kills principal ideles.
It kills norms because unramified quadratic norms have even
order, split norms are unrestricted, and complex norms are positive.
It is nontrivial on an idele negative at exactly one real place.
It therefore satisfies the literal \texttt{IsEta}.  The
\texttt{AllSome} clause in \texttt{ResEta} permits this particular
character.  At an inert place an $F_v$-uniformizer is also an
$E_P$-uniformizer, and consequently $\eta$ takes the value $-1$
on it.  It takes the value $1$ on all finite unit ideles by
unramifiedness.  This determines its entire inert factor.

Retain the same half selector $\varepsilon$ as in CW7.  For
$x\in E^\times$ let $\kappa(x)$ be its principal idele cut to
$\varepsilon$, and $\kappa'(x)$ the cut to
$\varepsilon+\bar\varepsilon$.
Write $\eta_{\rm in}(x)$ for the product over the inert places.
The Hecke property and the specified infinity type give
\[
 \eta(\kappa'(x))\eta_{\rm in}(x)
               \prod_\tau\frac{\tau(x)}{|\tau(x)|}=1.
                                                               \tag{CW.39}
\]
For integral $x\ne0$, factor its principal ideal.  At a fixed
prime $P$, CW2 gives $\lambda(NP)=-\sqrt{NP}$ and the preceding
local calculation gives $\eta_P(x)=(-1)^{\operatorname{ord}_P x}$.
At every moved prime the relevant rational primes are noninert,
so the corresponding $\lambda$ is the positive square root.
Multiplication proves
\[
 \lambda(|N_{E/\mathbb Q}x|)
       =\eta_{\rm in}(x)\sqrt{|N_{E/\mathbb Q}x|}.
\]
For arbitrary $x$ with $k=|N_{E/\mathbb Q}x|\in\mathbb N$, write
$x=u_0/z_0$ with $u_0,z_0$ nonzero algebraic integers.
The integer identity $k|N z_0|=|N u_0|$ and nonvanishing of
$\lambda(|Nz_0|)$ allow cancellation of the two integral formulas.
Combining with (CW.39) yields the exact positivity statement
\[
 \eta(\kappa'(x))\lambda(k)\prod_\tau\tau(x)
        =\sqrt{k}\prod_\tau|\tau(x)|>0.                \tag{CW.40}
\]
This uses the original $\eta,\lambda,\varepsilon$ and CM-type
embeddings; independent square-root choices would not give it.

\paragraph{CW19: explicit Gaussian theta tests and their zero coefficient.}
Fix Cartan involutions on $S,T$, and let $J$ be their real Hom
complex structure.  Write $\omega=\texttt{omO}$ and
$q(v)=\omega(v,Jv)>0$ for $v\ne0$.
Let $\iota$ be the inverse infinity chart and $\pi$ its projection.
For an adelic matrix $x$ define
\[
 \chi_A(x)=\mathbf1_A(x_f)\psi(b(wx_f,x_f)).
\]
For each real $a$ form the actual series
\[
 F_a(x)=\sum_{y\in V(E)}\psi(b(y,x))\,
    \psi(b(\iota a,y+x))\,\chi_A(y+x)
    \exp\bigl(-\tfrac\pi2 q(y_\infty+a+\pi x)\bigr).
                                                               \tag{CW.41}
\]
The rational vectors in this formula are diagonally embedded.
For fixed finite $x_f$, the support in $y$ is either empty or a
coset of $L$.  On a sufficiently small finite neighborhood the
same support and finite phases apply.  Real lattice-point counts
are polynomial, whereas the Gaussian and each derivative are a
polynomial times a positive Gaussian.  Therefore the series and
all real derivatives converge locally uniformly and absolutely.
This proves continuity and real smoothness, including exchange of
each derivative with the sum.

Reindexing $y$ proves the exact rational quasi-periodicity.
Expansion of the quadratic phase gives
\[
 \psi(-b(wt,t))\psi(b(x,t))\chi_A(x+t)=\chi_A(x)
                       \qquad(t\in A).
\]
It proves $W_tF_a=F_a$ term by term; real and finite components
have zero cross pairing.  On $2A$ the correcting phase is one,
so the raw finite Heisenberg condition holds on an open level.
The order $b(\iota a,y+x)$ in the second phase is essential:
it is the right magnetic translate of the Gaussian, which commutes
with the left Heisenberg derivatives; reversing its arguments
would not preserve the same vacuum equations.
Finally differentiating (CW.41) in directions $v,Jv$ cancels the
linear Gaussian derivative with the two Heisenberg phases, giving
$dH(v)F_a+i\,dH(Jv)F_a=0$.  Thus these are genuine theta functions
and simultaneous finite and Cartan vacua.

The family is not zero.  Setting $a=-v,x=\iota v$ expresses
$F_{-v}(\iota v)$ as an absolutely convergent Fourier series in
$v$, with frequency $y\mapsto-\omega(y_\infty,\cdot)$ and constant
coefficient $1$ from $y=0$.  All other frequencies are nonzero.
Self-duality makes them integral on $L$, so integration over its
real fundamental parallelepiped gives its positive volume.
In particular at least one $F_a$ is nonzero.  This supplies the
nonzero theta vector needed to make CW15's scalar unique.

Here is the precise trace calculation using these tests.  Use the
pair $p,u,C,n,i$ of CW17, and suppose that on every simultaneous
$A$-fixed Cartan vacuum one has
\[
 R_pf=\kappa\sum_{j=1}^{i}W_{b_j}f,\qquad
       b_j\in r(p)A/(A\cap r(p)A).                    \tag{CW.42}
\]
Define
\[
 \begin{split}
 Q_u(y)&=\mathbf1_{B_1}(y)\,
                \gamma_Q(\bar w(uy-y),y),\\
 G_u(v)&=e^{-\pi i\omega(uv,v)}
                     e^{-\pi q((u-1)v)/2}.
 \end{split}                                                   \tag{CW.43}
\]
Insert $F_{-v}$ in (CW.42) and evaluate at $\iota v$.
On the left, reindex by the bijection $y\mapsto(u-1)y$.
The phase equality is exactly
\[
 \chi_A((u-1)y)=Q_u(y)e^{-\pi i\omega(uy_\infty,y_\infty)},
\]
where the finite and real character values on a rational pairing
multiply to one.  Expanding the two remaining real phases then
gives $\sum_y Q_u(y)G_u(y_\infty+v)$.
On the right, the $j$th summand has Fourier coefficient
\[
 a_j(y)=\psi(-b(wb_j,b_j))\psi(b(y,b_j))
                    \chi_A(y+b_j)e^{-\pi q(y_\infty)/2}.
                                                               \tag{CW.44}
\]
These coefficients are absolutely summable over finitely many
lattice cosets.  The frequency is again
$-\omega(y_\infty,\cdot)$ and is integral on $C$ whenever
$a_j(y)\ne0$: test the self-dual lattices $A$ and $r(p)A$
against the finite images of $c$ and $uc$, for $c\in C$.
Only the zero frequency survives integration over a fundamental
parallelepiped for $C$.  Its total coefficient is
$\sum_j a_j(0)=1$, because exactly one representative belongs
to $A$, and there the two quadratic phases cancel.

On the left $Q_u$ is $C$-periodic, supported on the $n$ cosets
$t_k+C$.  Unfold each coset into the whole real vector space.
If $G_u$ is integrable, absolute integrability justifies this
operation and the right-hand Fourier integration.  CW20 proves
that integrability under the same condition $d\ne0$.
Writing $r_1>0$ for the real covolume of $C$, (CW.38) gives
\[
             \lambda(n)I_u=r_1\kappa,
                  \qquad I_u=\int G_u(v)\,dv.          \tag{CW.45}
\]
This proves the needed trace identity from the series, rather
than assuming it as a vacuum normalization axiom.

\paragraph{CW20: evaluating the real integral with its actual phase.}
Assume in addition that the rational pair commutes with the chosen
Cartans.  On the complex vectorization of the real Hom space let
$H$ be its original Hermitian matrix and $s$ its Cartan involution.
Thus $Hs$ is positive definite, $U^*HU=H$, and $Us=sU$;
the matrix $U$ here includes all CM-type embeddings of (CW.35).
The source convention is
$\omega(v_1,v_2)=-\Im(z_2^*Hz_1)$, so
$q(v)=z^*Hs z$.  Direct substitution in (CW.43) gives
$G_u(v)=e^{-\pi z^*E_u z}$ with
\[
 E_u=\tfrac12\bigl((U-1)^*Hs(U-1)-HU+(HU)^*\bigr).
                                                               \tag{CW.46}
\]
Its Hermitian part is
$\tfrac12(U-1)^*Hs(U-1)>0$, since $d\ne0$.
This proves absolute integrability.  CW9's complex Gaussian
formula therefore applies with no square-root ambiguity.
Put $P_+=(1+s)/2$, $P_-=(1-s)/2$.
Multiplying out (CW.46), using $U^*H=HU^{-1}$ and $Us=sU$,
gives the useful exact identity
\[
 E_u\,U(P_+U^{-1}+P_-)=-H(U-1).                       \tag{CW.47}
\]
Taking determinants and retaining the positive Jacobian between
the original real coordinates and these complex coordinates gives
\[
 I_u P_0=r_2\epsilon_0 Q_0\delta,\qquad r_2>0,\qquad
 P_0=\prod_\tau\tau(d),\quad Q_0=\prod_\tau\tau(\Delta),
                                                               \tag{CW.48}
\]
where $\epsilon_0=\det(-s)$ and
$\delta=\det(P_+U^{-1}+P_-)$.
Indeed $r_2$ is the coordinate Jacobian times $\det(Hs)^{-1}$,
which is positive real, and
$\det(-H)\det(-s)=\det(Hs)$ proves (CW.48).
The scalar $\delta$ is exactly CW13's real vacuum scalar for
$p$: $u$ was the inverse pair action.  It is nonzero by CW9,
and $\epsilon_0\ne0$ is fixed once the Cartans are fixed.
In particular no unknown complex factor is absorbed into $r_2$.

\paragraph{CW21: cancellation of all phases, including the finite index.}
Use the canonical characters
$\chi_S=\eta^{\dim S}$ and $\chi_T=\eta^{\dim T}$ in CW7.
For a rational pair commuting with the Cartans and satisfying
$d\ne0$, CW15, CW4 and CW13 give (CW.42) with
\[
       \kappa=c_p^{-1}\delta z\lambda(i)^{-1}.
                                                               \tag{CW.49}
\]
No uniqueness theorem for a merely unitary projective action is
being used: all three operators in this equation were constructed
on the same functions.  Thus CW19 gives
$\lambda(n)I_u=r_1c_p^{-1}\delta z\lambda(i)^{-1}$.
CW18 applied to the two numbers in (CW.37) gives positive reals
$r_3,r_4$ with
\[
 \begin{split}
 y_1\lambda(c_0)\lambda(n)ZP_0&=r_3,\\
 y_2\lambda(c_0)\lambda(i)ZQ_0&=r_4,
 \end{split}\quad
 \begin{gathered}
 y_1=\eta(\kappa'(N^m d)),\quad
 y_2=\eta(\kappa'(N^m\Delta)),\\
 Z=\prod_\tau\tau(N^m).
 \end{gathered}                                                \tag{CW.50}
\]
We check the remaining character identity explicitly.
Unitarity gives
\[
 \Delta\bar\Delta=1,\qquad
            \Delta\bar d=(-1)^m d.                    \tag{CW.51}
\]
The first is the determinant of $U^*HU=H$; the second follows
by taking determinants in
$U\bigl(\overline{U-1}\bigr)^{\mathsf t}$ after conjugation
by $H$, using $U^{-1}-1=-U^{-1}(U-1)$.
The cut identities and norm-triviality of $\eta$ give
\[
 \eta(\kappa'(x))\eta(\kappa(\bar x))=\eta(\kappa(x)),
 \qquad z\eta(\kappa(\Delta))=1.
\]
The last identity uses
$\Delta=\det(b_0)^{\dim S}\det(a)^{\dim T}$ and the
crossed character in (CW.11), including $b_0=h^{-1}$.
Applying these identities to (CW.51), and cancelling the common
$N^m$ factor, proves
\[
              zy_2=s_0y_1,\qquad
                 s_0=\eta(\kappa((-1)^m))\ne0.         \tag{CW.52}
\]
Inverting this order-two factor would give the same value;
no unspecified phase remains.

All factors cancelled below are nonzero: $d,\Delta,N$, the
character values, $c_p,\delta$, and positive-index values of
$\lambda$ have already been checked.  Multiply (CW.45) by
$c_p\lambda(i)$, substitute (CW.48), then eliminate $P_0,Q_0$
using (CW.50).  Finally use (CW.38) and (CW.52).  The result is
\[
             c_p\epsilon_0=r s_0,\qquad
                    r={r_1r_3\over r_2r_4 n}>0.       \tag{CW.53}
\]
This is where the same two rational lattice indices and the
signed Gauss sum are essential.  Absolute-value identities alone
would leave $c_p$ undetermined.

\paragraph{CW22: removing the no-fixed-vector restriction and proving automorphy.}
Keep any rational pair $p$ commuting with fixed Cartans, without
assuming $d\ne0$.  The norm-one elements of $E$ are infinite:
if $\xi\ne0$ is purely imaginary, then
$(a+\xi)/(a-\xi)$, $a\in\mathbb Q$, are distinct and have norm
one.  Choose such a $t$ outside the roots of
$\det(tU-1)$ and $\det(tI-1)$, two nonzero polynomials.
Let $k=tI$ on $T$ and $q=(1,k^{-1})$.
Then the inverse Hom actions for $q$ and $pq$ are $tI$ and $tU$.
They commute with the same Cartans and satisfy the no-fixed-vector
condition.  CW19 makes $\Theta$ nonzero, so the scalar in (CW.33)
is unique and the group law gives $c_{pq}=c_pc_q$.
Apply (CW.53) to $q$ and $pq$ and divide.  The same
$\epsilon_0,s_0$ cancel, proving $c_p$ is positive real.
Since $|c_p|=1$, necessarily $c_p=1$.
This argument also covers rank zero: the empty determinants and
zero-dimensional integrals are one.

Finally rational unitary groups are generated by the
quasi-reflections of WA10.  A quasi-reflection acts on a
nonisotropic line and fixes its orthogonal complement.
At each infinity embedding choose a Cartan decomposition preserving
that line, and choose positive and negative orthonormal bases on
its complement.  The resulting Cartan involution commutes with
the reflection.  For a reflection in one factor the identity in
the other factor commutes with any Cartan.  The preceding paragraph
therefore proves $c_p=1$ for each generator in either factor.
The actual group law extends it to all rational pairs.  We obtain
on the original theta space the full equality
\[
       \omega(\gamma,\gamma')f(x)
                =f\bigl(r(\gamma^{-1},\gamma'^{-1})x\bigr)
       \quad(\gamma\in G_S(E),\ \gamma'\in G_T(E)).     \tag{CW.54}
\]
Evaluating at zero gives rational theta-distribution automorphy
in both factors at once.  In particular this proves the literal
source and target automorphy clauses without a second representation
on transposed spaces.

\paragraph{CW23: the literal canonical full Weil datum.}
For the specified $D,\psi,S,T,\theta,\Lambda_S,\Lambda_T$ and
any supplied $\eta:\texttt{IsEtaTilde}$, the construction above
proves the actual conclusion of
\texttt{exists\_weilDatum\_etaTilde}.  We spell out its join.
The output is the group homomorphism on all functions obtained
by extending CW14's action on $\Theta$ once by an algebraic
complement.  Its characters are exactly
$\eta^{\dim S},\eta^{\dim T}$.  Powers preserve the Hecke
property, and raising the supplied restriction identity to a
power gives their two literal rank restrictions.

These are the first five fields of \texttt{IsWeilDatumO}, with
its standard-character field unchanged.  The remaining six are:
(CW4, CW13) preservation of the actual theta space;
(CW14) covariance with the full adelic Heisenberg action;
(CW22) rational evaluation-at-zero automorphy;
(CW14) pointwise continuity on each theta function;
(CW5, CW13) equality of the original Haar $L^2$ integrals;
and (CW6--CW7) the prescribed split factor times a positive
real scalar.  The norm equality holds for every additive Haar
measure and every rational fundamental domain in the definition,
so it also gives the literal nonnegative extended integral.

A split-supported group element has trivial real component:
apply the infinite projection to its support idempotent.
Thus its full action restricts to the finite action of CW7,
and the positive split-scalar clause is unchanged.  The same
argument for a nonsplit place shows that CW6 supplies
\texttt{FixesVacAt} for the original $\Lambda_S,\Lambda_T$.
At a real-supported pair the finite cut determinant is one and
its finite lattice action is the identity; hence CW13 supplies
\texttt{VacCharPlus}, with exactly (CW.29).
This proves both extra conclusions of the source theorem as
well as all eleven fields of its full datum.  It does not
replace them by a qualitative statement about algebraic jets.

\paragraph{CW24: prescribed splitting characters without an extra uniqueness premise.}
For arbitrary specified Hermitian spaces $S,T$ and specified
Hecke characters $\chi_S,\chi_T$ with the literal restrictions
$\texttt{ResEta}(\chi_S,\dim S)$ and
$\texttt{ResEta}(\chi_T,\dim T)$, a full oriented Weil datum
exists on these same spaces and functions.  We give a route
which also constructs the auxiliary character needed here,
without importing a new global norm-index theorem.

First use the actual standard lattice of each Hermitian space.
Its image under the given frame is
$\Lambda_S=\operatorname{frame}_S(S.L_0)$.
Finite generation and fullness are preserved by this linear
isomorphism.  The coordinate identity for the Hermitian form
shows that its dual is the image of the dual of $S.L_0$, hence
itself.  This proves the source's
\texttt{isSelfDualC\_latFrame}, and likewise for $T$;
no new Gram matrix is chosen.

HF11 supplies a unitary Hecke character $\eta_0$ of the actual
host, unramified at finite primes, with phase exponent one and
restriction equal to HF9's explicit quadratic character $\eta_F$.
Transport it to $(\mathbb A_F\otimes_F E)^\times$ by the
continuous tensor homomorphism
\[
 a\otimes x\longmapsto \operatorname{baseAdele}(a)
                                 \operatorname{diag}(x).
                                                               \tag{CW.55}
\]
On a finite component this is the completion extension, and at
infinity it is the specified embedding of $E$.  Agreement on
dense rational elements identifies each completion extension;
the finite-dimensional tensor presentation proves continuity.
It sends principal ideles to the same principal ideles, finite
integral units to finite integral units, and each local support
idempotent to its corresponding component.  Composing the HF11
character with (CW.55) therefore gives the same Hecke character
readings, infinity phases and inert values on the source's domain.
In particular its restriction is exactly $\eta_F$, and it kills
$x\bar x$ because $\eta_F$ kills norms.

There is a small logical distinction at this point.
\texttt{IsEtaTilde} contains a universal \texttt{AllSome} restriction
against every character satisfying \texttt{IsEta}.  Proving that
all those characters equal $\eta_F$ would be the separate global
norm-index-two uniqueness theorem.  We do not need that theorem.
Inspecting CW1--CW22, the construction of the operators uses only
that $\eta_0$ is Hecke, is unramified, has phase exponent one,
kills norms, and has the stated inert signs.  CW7 uses norm
triviality to define its cut character; CW18 uses precisely the
explicit local readings just proved.  Hence the same formulas
construct an action $\omega_0$ with covariance, continuity,
Haar norm, rational automorphy, and the positive split factor
for $\eta_0^{\dim S},\eta_0^{\dim T}$.  Only the two universal
restriction fields of a full record have not been asserted for
these auxiliary characters.

The supplied restrictions on $\chi_S,\chi_T$ can be tested
against the particular $\eta_F:\texttt{IsEta}$ from CW18.
Thus
\[
       \nu_S=\eta_0^{-\dim S}\chi_S,\qquad
       \nu_T=\eta_0^{-\dim T}\chi_T
\]
are Hecke characters trivial on \emph{all} ideles of $F$.
Apply the explicit determinant-root construction of WA18 to
these two characters and the action $\omega_0$:
\[
 \omega'(g,h)=\widetilde\nu_T(g)^{-1}
                            \widetilde\nu_S(h)\omega_0(g,h).
                                                               \tag{CW.56}
\]
WA18 proves root independence, restricted-product continuity,
rational triviality, and the exact crossed split-factor identity.
Those proofs use the operator properties just established, not
the missing universal restriction of $\eta_0$.
Consequently (CW.56) satisfies all operator clauses for the
prescribed $\chi_S,\chi_T$.  Its Hecke and two universal rank
restriction fields are the original supplied hypotheses.
This completes the full datum with no independent Weil-existence
or quadratic-character-uniqueness premise.  General character
twisting need not preserve the canonical extra vacuum determinant
or nonsplit fixing clauses; the full datum asserted here has
exactly its eleven stated fields.

\paragraph{CW25: actual bottom Weil families and the same bottom period field.}
Let $L$ be the supplied level-character datum with $L.\mathrm{Is}(\mu_0)$,
let $G$ be its supplied grid and $W:G.\mathrm{Weil}$ its supplied
realization.  For every square $(i,j)$, use WA17's actual standard
positive spaces $A,A',WW$ of ranks
$\mathrm{botBig},\mathrm{botSmall},\mathrm{botBig}+\mathrm{botSmall}$,
with their original identity frames and standard self-dual lattices.
Retain the prescribed bottom characters, including the exact
parity correction by $L.\eta^{\operatorname{rank}V\bmod2}$.
The rank identities in WA17 show that their restrictions have
exactly the required ranks.  CW24 now constructs, for every $t$,
the full representations on $(\operatorname{neg}A,WW)$ and
$(A',WW)$ with these characters and their prescribed character
quotients.  Negating the Hermitian form leaves the dual of the
same standard lattice unchanged.  The auxiliary $\eta_0$ of
CW24 does not replace $L.\eta$ in any bottom character.

This discharges the representation input explicitly retained in
WA17.  Its remaining construction chooses the normalized vacua,
finite idempotents, actual conductor exponents, standard finite
theta functions, Haar measures and fundamental domains on these
same spaces.  All identities refer to the same two
representations just constructed.  Thus the literal
\texttt{exists\_bottomTri} holds, including its actual
integrability clauses, by WA17 and TT12's definite quotient
argument.  No identification with the already supplied grid
representation $W$ is made.

If the twist set $T$ is finite, WA16 gives for each of these
actual bottom triangles $X_{ij}$ a finite number field
$M_{ij}\subset\mathbb C$ and nonzero periods for its two exact
integrals, satisfying \texttt{SumsIn}.  The compositum of the
finitely many $M_{ij}$ is finite.  Inclusion of fields preserves
all range-membership clauses without changing $X_{ij}$, its
measures or its periods.  This proves the literal bottom-field
portion \texttt{NFBot} for every supplied admissible realization;
its hypotheses provide this finiteness of $T$.

The canonical full bottom-Weil constructor is therefore no
longer a retained primitive.  The quantitative all-embedding
finite-unit dictionary estimates, the critical toric ratio field,
and the remaining same-choice \texttt{Choice.FieldOK} embedding
join retain the separate scopes recorded in WA19.  None is a
consequence of (CW.54) alone, and the supplied arithmetic host,
level data and grid realization have not been constructed anew.
